\documentclass[11pt, a4paper]{article}
\usepackage[utf8]{inputenc}
\usepackage[T1]{fontenc}
\usepackage[ngerman]{babel}
\usepackage{geometry}
\usepackage{helvet}
\usepackage{parskip}
\usepackage{csquotes} % Für saubere Anführungszeichen

% Seitenränder anpassen
\geometry{left=2.5cm, right=2.5cm, top=2.5cm, bottom=2.5cm}

% Schriftartwechsel auf Sans-Serif
\renewcommand{\familydefault}{\sfdefault}

\begin{document}

\section*{One-Pager / Kurzzusammenfassung}
\textbf{Arbeitstitel:} Konzeption einer Informationsverarbeitungsarchitektur für \enquote{Architecture as Code}: Bewertungs- und Synthesestrategien zur sicheren Überführung von Bestandsdaten in SysML v2

\textbf{Bearbeiter:} Moritz Diez \quad | \quad \textbf{Datum:} \today

\vspace{0.5cm}
\hrule
\vspace{0.5cm}

\textbf{Problemstellung:} \\
Im industriellen Umfeld liegen Informationen zu Bestandsystemen häufig unstrukturiert und heterogen vor (\enquote{Schuhschachtel-Dilemma}). Der Einsatz generativer KI bietet zwar neue Möglichkeiten zur Automatisierung, birgt jedoch ohne klare Architektur erhebliche Risiken: Wenn KI-Agenten autonom über die Relevanz und Interpretation unstrukturierter Daten entscheiden, drohen Halluzinationen und \enquote{Garbage In, Garbage Out}-Effekte. Es fehlt eine Referenzarchitektur, die definiert, welche Qualitätsanforderungen (Constraints) Eingangsdaten erfüllen müssen und wie der Ingenieur effektiv als Kontrollorgan eingebunden wird, um die technische Korrektheit des Zielmodells sicherzustellen.

\vspace{0.5cm}

\textbf{Lösungsansatz:} \\
Die Arbeit konzipiert eine Informationsverarbeitungsarchitektur, die SysML v2 als Enabler nutzt, aber den Fokus auf die Prozess-Sicherheit legt. Anstelle einer reinen \enquote{Knopfdruck-Automatisierung} wird ein assistierter Prozess (\enquote{Human-in-the-Loop}) entwickelt:
\begin{enumerate}
    \item \textbf{Data Assessment Framework:} Eingangsdaten werden vor der Verarbeitung gegen definierte Constraints geprüft (z. B. Maschinenlesbarkeit). Daten, die für eine automatisierte Verarbeitung zu mehrdeutig sind, werden frühzeitig erkannt.
    \item \textbf{Smart Agent \& Gatekeeping:} Ein KI-Agent antizipiert fehlende Zusammenhänge (\enquote{Pending Interfaces}), fordert jedoch bei Ontologie-Konflikten explizit eine menschliche Entscheidung an (\enquote{Effective Intervention}).
    \item \textbf{Ontology-Guided Synthesis:} Die Codegenerierung erfolgt innerhalb starrer ontologischer Leitplanken, um validen und modularen SysML v2 Code zu gewährleisten.
\end{enumerate}

\vspace{0.5cm}

\textbf{Geplante Validierung:} \\
Die Tragfähigkeit der Architektur wird anhand eines \enquote{Vertical Slice} (einem durchgängigen Pfad von der Rohdaten-Bewertung bis zum Code) validiert. Die Evaluation erfolgt nicht gegen eine Musterlösung, sondern anhand von Architektur-Metriken in Anlehnung an ISO 42030 (z. B. Gatekeeping-Effizienz, Ontologie-Treue). Damit wird aufgezeigt, dass erfolgreiches \enquote{Architecture as Code} primär durch intelligentes Architektur-Design und Daten-Assessment gelingt, nicht allein durch KI-Leistung.

\end{document}